\documentclass[aspectratio=169,11pt]{beamer}

% Packages (Matching final_presentation.tex)
\usepackage[utf8]{inputenc}
\usepackage[T1]{fontenc}
\usepackage{lmodern}
\usepackage{microtype}
\usepackage{amsmath,amssymb}
\usepackage{booktabs}
\usepackage{graphicx}
\usepackage{tikz}
\usepackage{pgfplots}
\usepackage{xcolor}
\usepackage{ragged2e}
\usepackage{colortbl}
\usepackage{array}
\usepackage{hyperref}
\usepackage{adjustbox}

\usetikzlibrary{shapes.geometric, arrows.meta, positioning, calc,
                backgrounds, decorations.pathreplacing, shadows, fadings}
\pgfplotsset{compat=1.18}

% --- Color Palette: Warm Professional ---
\definecolor{DeepNavy}{HTML}{2E4057}
\definecolor{Teal}{HTML}{048A81}
\definecolor{WarmOrange}{HTML}{E85D04}
\definecolor{SoftPurple}{HTML}{9D4EDD}
\definecolor{WarmGray}{HTML}{6C757D}
\definecolor{LightGray}{HTML}{E9ECEF}
\definecolor{Cream}{HTML}{FBF8F1}
\definecolor{DeepRed}{HTML}{D62828}
\definecolor{Gold}{HTML}{D4A03A}
\definecolor{SoftWhite}{HTML}{FAFAFA}

% Beamer colors
\setbeamercolor{normal text}{fg=DeepNavy, bg=SoftWhite}
\setbeamercolor{structure}{fg=DeepNavy}
\setbeamercolor{alerted text}{fg=DeepRed}
\setbeamercolor{example text}{fg=Teal}
\setbeamercolor{frametitle}{fg=DeepNavy, bg=SoftWhite}
\setbeamercolor{title}{fg=DeepNavy}
\setbeamercolor{subtitle}{fg=WarmGray}
\setbeamercolor{block title}{bg=DeepNavy, fg=SoftWhite}
\setbeamercolor{block body}{bg=LightGray, fg=DeepNavy}
\setbeamercolor{itemize item}{fg=WarmOrange}
\setbeamercolor{itemize subitem}{fg=Gold}

% Fonts
\usefonttheme{professionalfonts}
\setbeamerfont{title}{size=\huge, series=\bfseries}
\setbeamerfont{frametitle}{size=\Large, series=\bfseries}

% Minimal Templates
\setbeamertemplate{navigation symbols}{}
\setbeamertemplate{headline}{}
\setbeamertemplate{footline}{%
    \hfill
    \begin{beamercolorbox}[wd=3cm,ht=2.5ex,dp=1ex,right,rightskip=0.6cm]{page number}
        \usebeamercolor[fg]{WarmGray}\scriptsize Supplemental Slide
    \end{beamercolorbox}
    \vspace{0.3cm}
}

\begin{document}

% --- SLIDE 1: THE PRACTICAL CHALLENGE ---
\begin{frame}
    \frametitle{Latent macro-sensitivity creates "leaky" instruments}
    
    \vspace{0.5cm}
    {\Large
        In practice, we rarely know the true factor loadings ($\lambda_i$)—how each firm responds to the macroeconomy.
    }
    
    \vspace{0.8cm}
    \begin{itemize}
        \item \textbf{The Risk:} If we assume $\lambda_i = 1$ (Simple GIV) when large firms are systematically more sensitive to macro shocks, the instrument becomes correlated with $\eta_t$.
        \item \textbf{The Consequence:} The exclusion restriction is violated, and IV estimates of elasticities are biased.
        \item \textbf{The Solution:} We must estimate the loadings ($\hat{\lambda}_i$) from the panel data to construct the orthogonal weights $\Gamma$.
    \end{itemize}
\end{frame}

% --- SLIDE 2: PCA METHOD ---
\begin{frame}
    \frametitle{PCA extracts factor loadings from firm-level outcomes}
    
    \vspace{0.5cm}
    {\Large
        If we assume idiosyncratic shocks are small relative to common factors, we can use \textbf{Principal Component Analysis}.
    }
    
    \vspace{0.8cm}
    \begin{enumerate}
        \item \textbf{Panel Analysis:} Apply PCA to the $T \times N$ matrix of outcomes ($L_{it}$).
        \item \textbf{Factor Extraction:} The first $K$ principal components represent the latent macro shocks ($\hat{\eta}_t$).
        \item \textbf{Loadings:} The coefficients of the firm outcomes on these components are the estimated loadings ($\hat{\lambda}_i$).
        \item \textbf{Weighting:} Construct $\Gamma$ by projecting firm sizes $S$ orthogonal to $\hat{\lambda}$ and a constant.
    \end{enumerate}
\end{frame}

% --- SLIDE 3: GMM METHOD ---
\begin{frame}
    \frametitle{Iterative GMM identifies parameters and loadings simultaneously}
    
    \vspace{0.3cm}
    {\Large
        For a more robust approach, Gabaix and Koijen (2022) suggest joint estimation via the \textbf{Generalized Method of Moments}.
    }
    
    \vspace{0.5cm}
    \textbf{The Iterative Logic:}
    \begin{itemize}
        \item \textbf{Step 1:} Start with an initial guess for the elasticity $\psi^{(0)}$.
        \item \textbf{Step 2:} Extract the "macro-contaminated" residuals from the demand equation.
        \item \textbf{Step 3:} Estimate $\lambda_i$ from these residuals (e.g., via the first principal component).
        \item \textbf{Step 4:} Update the GIV instrument $z_t$ using the new loadings and re-estimate $\psi^{(1)}$.
    \end{itemize}
    \vspace{0.3cm}
    \textit{This process continues until the elasticity and factor loadings converge.}
\end{frame}

% --- SLIDE 4: SUMMARY TABLE ---
\begin{frame}
    \frametitle{Choosing the estimation strategy in practice}
    
    \vspace{0.5cm}
    \begin{center}
    \renewcommand{\arraystretch}{1.5}
    \begin{tabular}{l l l}
        \toprule
        \textbf{Assumption} & \textbf{Approach} & \textbf{Estimation Method} \\
        \midrule
        \rowcolor{Teal!10} Homogeneous Loadings & Simple GIV & 2SLS (Difference-in-Averages) \\
        Known Factor Loadings & Optimal GIV & 2SLS with Analytical $\Gamma$ \\
        \rowcolor{WarmOrange!10} \textbf{Unknown Loadings} & \textbf{Latent GIV} & \textbf{PCA or Iterative GMM} \\
        \bottomrule
    \end{tabular}
    \end{center}
    
    \vspace{0.8cm}
    \textbf{Practical Note:} When using estimated loadings ($\hat{\lambda}$), standard errors should ideally be bootstrapped or adjusted to account for the "generated regressor" variance.
\end{frame}

\end{document}
